\begin{table}[htbp]
\centering
\caption{Observed Attack Composition and Repeat Metrics by Trade-Size Tier}
\label{tab:victim_tiers}
\begin{tabular}{lrrrrrrrrr}
\toprule
Tier & Victim Trades (N) & Tier Address Obs (q5b) & Trader EOAs (tx\_from, non-bot) & Trader Attacks (N) & Attacks / Trader & Median Attacks / Trader & Known-Bot EOAs & Known-Bot Attacks / EOA & Attacks / $1k Traded \\
\midrule
Retail & 1,876,928 & 209,181 & 393,766 & 1,873,714 & 4.7584 & 1 & 684 & 92.1579 & 2.4400 \\
Small & 1,501,543 & 142,031 & 278,906 & 1,507,105 & 5.4036 & 1 & 205 & 187.7854 & 0.3244 \\
Institutional & 375,386 & 40,517 & 68,732 & 267,152 & 3.8869 & 1 & 37 & 117.6757 & 0.0170 \\
\bottomrule
\end{tabular}
\vspace{0.2cm}
\begin{minipage}{\linewidth}
\small
\textit{Notes:} This table reports observed attacked-trade composition by trade-size tier (query5b_victim_impact_v2) and tx_from identity-corrected repeat metrics (query5c_repeat_victimization_v2) across 3,753,857 attack events. Tier-level address counts in q5b (391,729) are non-additive tier-address observations and must not be interpreted as globally unique victims. Primary identity counts are tx_from trader EOAs excluding known bots (741,404); known-bot victim EOAs (926) are reported separately. \textbf{Median finding:} the median attacks per trader EOA equals 1 in every tier (Retail, Small, Institutional), so the repeat-victimisation hypothesis is withdrawn as a typical-case claim. \textbf{Mechanical identity caveat:} Attacks per $1,000 traded equals 1,000 divided by average trade size by construction. Accordingly, the Retail/Institutional attacks-per-$1k ratio (143.82x) is algebraically identical to the Institutional/Retail mean-size ratio (143.82x).
\end{minipage}
\end{table}
